\subsection{记号与术语}

本节其余部分将要给出的定义适用于x和y之间的任意序关系（或预序关系）R\{x, y\}，但主要将在R\{x, y\} 被写成 \(x \leqslant y\) *（类比整数或实数之间通常的序关系）*（或 \(x \subset y\) ，或通过某种类似的符号）的情况下使用；因此我们仅针对记号 \(x \leqslant y\) 来陈述它们，而将把它们转写为其他记号的任务留给读者。当 R\{x, y\} 被写成 \(x \leqslant y\) ，那么 \(y \geqslant x\) 与 \(x \leqslant y\) 同义，这些关系读作“x小于y”，或“x比y小”，或“y大于x”，或“y比x大”。关系 \(x \geqslant y\) 于是是与 \(x \leqslant y\) 相反的（在x与y之间的）预序关系。

滥用语言地，我们常常会说“关系 \(\leqslant\) '' 而不是 `` 关系 \(x \leqslant y\) ''；在这种情况下，`` 关系 \(\geqslant\) '' 是 `` 关系 \(\leqslant\) '' 的相反关系。我们还注意到，在同一证明中，我们常常可以使用相同的符号 \(\leqslant\) 用于在不会引起混淆的情况下表示几种不同的序关系。

写作 \(x \leqslant y\) 的关系是集合 \(\mathbf{E}\) 上的一个序关系，其条件如下：

\((\mathbf{RO}_{\mathbf{I}})\) 关系“ \(x \leqslant y\) 并且 \(y \leqslant z\) '' 蕴含 \(x \leqslant z\) .

\((\mathbf{RO}_{\mathbf{II}})\) 关系“ \(x \leqslant y\) 并且 \(y \leqslant x\) '' 蕴含 \(x = y\) .

\((\mathbf{RO}_{\mathbf{III}})\) 关系 \(x \leqslant y\) 蕴含“ \(x \leqslant x\) 并且 \(y \leqslant y\) ''.

\((\mathbf{RO}_{\mathbf{IV}})\) 关系 \(x \leqslant x\) 等价于 \(x \in \mathbf{E}\) .

如果我们省略条件 \((\mathrm{RO}_{\mathbf{II}})\) ，我们就得到 \(x \leqslant y\) 是 E 上的预序关系的条件。

当序关系写作 \(x \leqslant y\) 时，我们将用 x \textless{} y（或 y \textgreater{} x）表示关系“ \(x \leqslant y\) 并且 \(x \neq y\) ”；这些关系读作“x 严格小于 y”，或“x 比 y 小”，或“y 严格大于 x”，或“y 比 x 大”。

包含关系的例子表明， \(x \leqslant y\) 的否定（有时记作 \(x \nleqslant y\) ）不一定等价于 \(y < x\) （参见第 12 号）。

C58. 设 \(\leqslant\) 是一个序关系，并设 x, y 是两个不同的字母。关系 \(x \leqslant y\) 等价于“x \textless{} y 或 x = y”。关系“ \(x \leqslant y\) 并且 \(y < z\) '', `` \(x < y\) 并且 \(y \leqslant z\) ”中的每一个都蕴含 x \textless{} z。

第一个断言由准则 \(\mathbf{A} \Longrightarrow ((\mathbf{A}\) 以及（非 \(\mathbf{B}\) ）或 \(\mathbf{B}\) ）（第一章，§3，判据C24）。为证明第二个断言，我们注意到每个假设都蕴含 \(x \leqslant z\) ，由传递性；而关系（ \(x = z\) 并且 \(x \leqslant y\) 并且 \(y \leqslant z\) ）将蕴含 \(x = y = z\) ，与假设矛盾。

为使问题更简便，并用数学定理取代元数学判据，我们通常考虑一个理论 \(\mathfrak{G}\) ，它包含集合论的公理及公理模式，此外还有两个常项 \(\mathbf{E}\) 与 \(\Gamma\) 满足公理

\textbf{“Γ是集合E上的一个序关系”（第1节）。}

我们将用 \(x \leqslant y\) 表示关系 \(y \in \Gamma \langle x \rangle\) ，并且我们说集合 \(E\) 由序 \(\Gamma\) 排序（或由序关系 \(y \in \Gamma \langle x \rangle\) ）（参见第四章，§1）。

¶ 每当在 \(\mathfrak{G}\) 中 \(\Gamma\) 是 \(\mathbf{E}\) 上的预序，我们同样称 \(\mathbf{E}\) 由预序 \(\Gamma\) 所预序.

在某些情形下（例如在下面的定义中），我们将要考虑的理论稍微复杂一些。我们将把这些理论的常项和公理明确化的工作留给读者。

设 E， \(E'\) 为两个由序 \(\Gamma\) , \(\Gamma'\) 排序的集合。从 E 到 \(E'\) （对于序 \(\Gamma\) 与 \(\Gamma'\) ）的一个同构是 E 到 \(E'\) 上的一个双射 f，使得关系 \(x \leqslant y\) 与 \(f(x) \leqslant f(y)\) 等价（参见第四章，§1）。

